\section{Síntesis y ampliación}
\begin{tabular}{@{}p{4.5cm}p{11cm}@{}}
\toprule
Dimensión & Puntos clave \\
\midrule
Panorama del escenario & ServiceNow, mediante la historia de John/Sandy, revela oportunidades de automatización de tickets poco frecuentes pero críticas. \\
Ruta de arquitectura & El API Agent usa un conjunto de herramientas normalizado, el Web Agent opera el navegador; ambos deben complementarse. \\
TapeAgents & El tape es un registro estructurado que permite depurar, optimizar y destilar, y resuelve la separación entre el desarrollo del prompt y el del modelo. \\
Ciclo de optimización & Se usa un modelo grande para hacer demostraciones y el tape para producir datos con los que se hace fine-tuning de un modelo pequeño, lo que hace que la calidad de GRATE sea visible y el costo, controlable. \\
Requisitos de gobernanza & Auditoría de extremo a extremo, aprobación de cumplimiento normativo; el prompt y el resumen de ejecución deben poder consultarse. \\
\bottomrule
\end{tabular}

\subsection{Lecturas adicionales}
\begin{itemize}
    \item El informe técnico oficial de TapeAgents y los recursos de GitHub (el presentador da explícitamente el código QR y el enlace en 00:19:05).\newline
    \item El artículo original de WorkArena/BrowserGym: ofrece una evaluación estandarizada y un benchmark para Web Agents de nivel empresarial.
    \item La presentación de la «Enterprise Workflow Platform» de ServiceNow ayuda a entender las prácticas estándar de colaboración entre múltiples sistemas y de gobernanza.
\end{itemize}

\end{document}
